% Presentación abreviada y precisa del contenido de un documento, sin agregar interpretación o crítica; no debe exceder una página.
\section*{Resumen}

\pendiente{Cifras provisionales: juez de relevancia local, conjunto de consultas en borrador y validación humana del juez pendiente. Se actualizan con la corrida final.}

La búsqueda de empleo remoto se reparte entre varias plataformas cuyos motores comparan las palabras de la consulta con las de la oferta, y fallan cuando ambas describen el mismo cargo con términos distintos. Este trabajo desarrolla y evalúa JobAI, una plataforma de búsqueda de empleo remoto basada en un pipeline de generación aumentada por recuperación (RAG), con el objetivo de superar las limitaciones de precisión de la búsqueda por palabras clave. Siguiendo el paradigma de Design Science Research, el trabajo constituye una nueva iteración del ciclo de diseño-evaluación del artefacto. En la fase de diseño se corrigió la consulta vectorial, que no usaba el índice, y se migró el almacenamiento vectorial de pgvector a Qdrant detrás de una interfaz común. La evaluación se realizó sobre un corpus congelado de 1.068 ofertas reales (4.013 fragmentos) y 100 consultas: 75 en inglés, en cinco categorías, y 25 en español. La recuperación densa, servida por pgvector y por Qdrant, se comparó con la búsqueda por subcadena que ofrece la plataforma (ILIKE) y con BM25. Los juicios de relevancia los asignó un modelo de lenguaje sobre un agrupamiento de profundidad 10. La fidelidad de las respuestas se midió con RAGAS, la latencia por etapa del pipeline y el compromiso entre \textit{recall} y latencia de ambos motores sobre el corpus real y conjuntos sintéticos de 100.000 y 500.000 vectores, con y sin filtros.

La recuperación densa obtuvo un nDCG@10 de 0,73, frente a 0,56 de BM25 y 0,06 de ILIKE. Las diferencias fueron significativas en las cinco métricas según la prueba de aleatorización pareada con corrección de Holm. La mayor ventaja sobre BM25 se dio en las consultas con desajuste de vocabulario (+0,21). En español, el nDCG@10 de la recuperación densa estuvo entre 0,64 y 0,66, y el de BM25 fue de 0,34. La búsqueda aproximada con HNSW no perdió calidad frente a la búsqueda exacta. En las consultas de agregación, el clasificador eligió la ruta y la plantilla SQL esperadas en 12 de 15 casos, y la respuesta coincidió con la consulta de referencia en 9. Con el juez local, la fidelidad media fue de 0,663 y la relevancia de la respuesta de 0,599. Sin embargo, la revisión de los casos con puntaje nulo mostró errores del juez, por lo que estas cifras no se consideran una medida fiable. Servido desde Qdrant, el contexto se construyó en 59~ms en la mediana, frente a 80,5~ms con pgvector, cuyo planificador recorre la tabla completa a esta escala. La generación dominó la latencia: el primer \textit{token} llegó en unos 0,53~s y la respuesta completa en unos 1,21~s. En el benchmark sobre el corpus real y sin filtro, pgvector igualó o superó el \textit{recall} de Qdrant con menor latencia hasta $ef = 128$. Con un filtro del 50\,\% y sin barrido iterativo, pgvector devolvió los 10 resultados solo en el 25,7\,\% de las consultas, mientras que Qdrant mantuvo un \textit{recall} de 1,000 con todos los filtros. \pendiente{Resultados con 100.000 y 500.000 vectores.}

\textbf{Palabras clave}: generación aumentada por recuperación, búsqueda semántica, búsqueda de empleo, bases de datos vectoriales, pgvector, Qdrant, evaluación de la recuperación de información, Design Science Research.

\newpage
\section*{Abstract}

\pendiente{English version of the abstract.}

\textbf{Keywords}: \pendiente{keywords}.
